%% c27 —— 印刷页 223–235（PDF 224–236）
%% 原文：G. W. Moore, Introduction To Chern-Simons Theories (TASI 2019)

\subsubsection{通往 Chern-Simons 的其他路径} % 原文: I.7.1 Other paths to Chern-Simons

对 Chern-Simons 形式还有一种更为概念化的描述，如下所述。考虑联络的一条任意路径 $\gamma$，即 $A(t)$。利用从 $M\times I\to M$ 的拉回丛，我们得到丛 $\pi^*(P)\to M\times I$，而这一族联络定义了一个联络：
\begin{equation}
d_M+dt\,\frac{\partial}{\partial t}+A(t)
\tag{I.55}
\label{eq:I.55}
\end{equation}
它定义在 $\pi^*(P)$ 上。在 $\Omega^{p,q}(M\times I)$ 上引入双分次（bigrading）是有用的，使得 $(p,q)$ 型形式落在 $\Omega^p(M)\otimes\Omega^q(I)$ 之中。局部地，我们的联络形式 $A(t)$ 是 $(1,0)$ 型的。让我们把对应于协变导数 \eqref{eq:I.55} 的联络记作 $\mathcal{A}$。由于 $A(t)$ 依赖于 $t$，曲率 $F(\mathcal{A})$ 同时具有 $(2,0)$ 型分量与 $(1,1)$ 型分量：
\begin{equation}
F(\mathcal{A})=F(A(t))+dt\wedge\frac{\partial A}{\partial t}
\tag{I.56}
\label{eq:I.56}
\end{equation}
（注意 $\frac{\partial A}{\partial t}\in\Omega^1(M;\operatorname{ad}P)$。）现在我们可以定义一个更一般的 Chern-Simons 形式：
\begin{equation}
CS_\gamma:=\int_I\operatorname{Tr}F(\mathcal{A})^n
\tag{I.57}
\label{eq:I.57}
\end{equation}
对任意路径 $\gamma$ 都成立。若 $\gamma(0)=A_0$、$\gamma(1)=A_1$，则它满足
\begin{equation}
d_M CS_\gamma=\operatorname{Tr}F(A_1)^n-\operatorname{Tr}F(A_0)^n
\tag{I.58}
\label{eq:I.58}
\end{equation}
上面这些 Chern-Simons 形式是通过在 $A_0$ 与 $A_1$ 之间选取线性路径而定义出来的。

% [Tutor Note: (I.56) 的类型分解依赖于双分次约定：$A(t)$ 是 $(1,0)$ 型，$dt$ 是 $(0,1)$ 型，而全外微分应理解为 $d=d_M+dt\,\partial_t$。由于 $dt$ 要与 $A$ 反交换才能写成 $dt\wedge\partial_tA$，若把 $\mathcal{A}$ 写成"$1$ 形式" $A(t)+dt\,\partial_t$ 直接代入 $F=d\mathcal{A}+\mathcal{A}\wedge\mathcal{A}$，$(1,1)$ 分量的符号取决于 $\partial_t$ 是作用在系数上还是被当作沿 $I$ 的额外分量；核对时应回到 $d\mathcal{A}=d_MA+dt\wedge\partial_tA$ 这一行，并注意 $\partial_tA\in\Omega^1(M;\operatorname{ad}P)$ 仍是 $M$ 上的 $1$ 形式。]

% [Tutor Note: (I.57) 里 $\int_I$ 的合法性藏了一步：$\operatorname{Tr}F(\mathcal{A})^n$ 按双分次展开为 $\operatorname{Tr}F(A(t))^n+n\,dt\wedge\operatorname{Tr}\big(\partial_tA\,F^{n-1}\big)$（各项的 $dt$ 次数不同），只有含 $dt$ 恰好一次的那一项在 $\int_I$ 后存活，因此 $CS_\gamma$ 是 $M$ 上的 $(2n-1)$ 形式；作者未写出该投影。把它代回 (I.58) 时应看到 $d_M$ 只作用在 $M$ 的方向上，$dt$ 项正是 $\partial_t\operatorname{Tr}F(A(t))^n$ 的全导数，这才是 (I.58) 右端出现 $A_1$ 与 $A_0$ 之差的原因。]

一个能说明「路径选取确实要紧」的例子如下：考虑一族规范变换 $g(t)$，$0\le t\le 1$，它定义了从 $A$ 到 $A^g$ 的一条路径，其中 $g=g(1)$。注意这条路径不同于线性路径。

假设我们有两条路径 $A_0+\alpha_1(t)$ 与 $A_0+\alpha_2(t)$。由于联络空间是仿射空间，我们可以找到这两条路径之间的同伦 $\alpha(s,t)$。

于是：$\alpha(0,t)=\alpha_1(t)$，$\alpha(1,t)=\alpha_2(t)$，且 $\alpha_1(1)=\alpha_2(1)=\alpha$。那么
\begin{equation}
CS_{\gamma_1}-CS_{\gamma_2}=d\int_{I\times I}\operatorname{Tr}F(\mathcal{A})^n
\tag{I.59}
\label{eq:I.59}
\end{equation}
其中对 $I\times I$ 的积分是对无穷维向量空间 $\Gamma(\Omega(M;\operatorname{ad}P))$ 中张成这两条路径的某个曲面所作的积分，其定向取得恰当。

% [Tutor Note: (I.59) 的两处跳步：(i) 右端的 $d$ 是 $M$ 上的外微分而非 $M\times I$ 上的全微分，故只在把 $\alpha(s,t)$ 视为 $I\times I\to\Gamma(\Omega(M;\operatorname{ad}P))$ 的映射并取 $(n,n)$ 分量之后才有意义；(ii) 这只说明 $CS_{\gamma_1}-CS_{\gamma_2}$ 是恰当形式，其前提是同伦 $\alpha(s,t)$ 光滑且两端 $\alpha_1(1)=\alpha_2(1)=\alpha$ 固定——若两端点不固定，差就不再恰当而是差一个 $\operatorname{Tr}F^n$ 的积分。原文"suitably oriented"一句即指 $I\times I$ 的定向需与 $\gamma_1,\gamma_2$ 的次序相容，否则整体符号翻转。]

{\small\itshape ［原文标记：$\ast\ast\ast\ast\ast$ 仍需完成：1. $\int_{M^3}CS(A)$ 作为 $\mathbb{R}/2\pi\mathbb{Z}$ 中元素之取值；2. Lie 代数上不变形式的仔细归一化；3. Chern-Simons 作为二次精化（quadratic refinement）：归入微分上同调（diff'l coho）一节。$\ast\ast\ast\ast\ast$］}
% [原文待办标记，原文 p.223；原文以两行星号行围出 STILL NEEDED 清单，逐项照译]

\section{Lie 群与 Lie 代数的一些材料} % 原文: J. Some Material On Lie Groups And Lie Algebras

\begin{keytakeaways}
  \item 结构定理：$G$ 的单位连通分支 $G_1$ 具有形式 $G_1\cong(G^{ss}\times U(1)^r)/Z$，其中 $G^{ss}$ 是紧、连通、半单且单连通的 Lie 群，而 $Z$ 是落在 $G^{ss}\times U(1)^r$ 中心内的有限阿贝尔子群。
  \item 整体拓扑：$G/G_1$ 恰是连通分支群，故有扩张 $1\to G_1\to G\to\pi_0(G)\to 1$；该扩张一般并不分裂，$O(2)=SO(2)\rtimes\mathbb{Z}_2$ 就是反例，而任何 Lie 群的 $\pi_1$ 恒为阿贝尔群。
  \item 单连通单群的清单（$A_r$ 至 $G_2$）连同其中心 $Z(\mathfrak{g})$ 与外自同构群，决定了 $G^{ss}$ 的全部可能整体形式；正维数群的扩张存在时自动给出典范同态 $\bar\omega:\pi_0(G)\to\operatorname{Out}(G_1)$。
  \item 对偶 Coxeter 数由小表给出：$A_r\mapsto r+1$，$B_r\mapsto 2r-1$，$C_r\mapsto r+1$，$D_r\mapsto 2r-2$，$E_8\mapsto 30$，$G_2\mapsto 4$。脚注特别警告：对非单连通（nonsimply laced）的代数另有一个 Coxeter 数，与对偶 Coxeter 数并不相等。
\end{keytakeaways}

\subsection{一般紧 Lie 群的结构} % 原文: J.1 The Structure Of The General Compact Lie Group

若 $G$ 是 Lie 群，则其单位连通分支（记作 $G_1$）是正规子群。\footnote{证明：若 $g\in G_1$，取连接 $g$ 与 $1_G$ 的路径 $g(t)$。若 $h\in G$，则 $h^{-1}g(t)h$ 连接 $h^{-1}gh$ 与 $1_G$。} 因此 $G/G_1$ 是一个群，但它不过是 $G$ 的连通分支群，于是我们有
\begin{equation}
1\to G_1\to G\to\pi_0(G)\to 1
\tag{J.1}
\label{eq:J.1}
\end{equation}
一般地 $G$ 并不是直积：作为群，$G$ 不同构于 $G_1\times\pi_0(G)$。举个简单的例子：$O(2)=SO(2)\rtimes\mathbb{Z}_2$。作为流形它是两个圆的乘积，但 $\mathbb{Z}_2$ 因子通过非平凡自同构 $R(\theta)\to R(-\theta)$ 而作用。

{\small\itshape ［原文标记：给出一个它不分裂的例子］}
% [原文边注/待办标记，原文 p.224；原文以 \ensuremath{\clubsuit}…\ensuremath{\clubsuit} 标在页边，内容是 "Give example where it does not split"]

现在，$\pi_0(G)$ 可以是任意的有限群（毕竟我们可以干脆取 $G$ 为该有限群！）。

现在我们集中考察 $G_1$。它是连通的紧 Lie 群。其万有覆盖 $\widetilde{G_1}$ 也是连通的 Lie 群，尽管可能是非紧的。可以证明，任何 Lie 群的 $\pi_1$ 都是阿贝尔群\footnote{证明：给定两条路径 $g_1(t),g_2(t)$，考虑连续函数 $F(t,s)=g_1(t)g_2(s)$。}，因此（假定 $\pi_1$ 有限生成）它非典范地是一个有限群与格 $\mathbb{Z}^d$ 的乘积。该格必须作用在 $\mathbb{R}^d$ 上以给出 $U(1)^d$，「于是」（我们这里并没有给出完全严格的论证，但结论是对的）$G_1$ 具有如下形式
\begin{equation}
G_1\cong(G^{ss}\times U(1)^r)/Z
\tag{J.2}
\label{eq:J.2}
\end{equation}
其中 $G^{ss}$ 是紧、连通、半单且单连通的 Lie 群，而 $Z\subset G^{ss}\times U(1)^r$ 是中心的一个子群，并且是有限阿贝尔群。进而 $G^{ss}$ 具有形式
\begin{equation}
G^{ss}=\prod_i G^{ss}_i
\tag{J.3}
\label{eq:J.3}
\end{equation}
其中 $G^{ss}_i$ 是紧、连通、单且单连通的 Lie 群。这些群的清单，连同它们的中心与外自同构群，给出在下面的表中。

\begin{center}
\begin{tabular}{|c|c|c|}
\hline
$\mathfrak{g}$ & $\widetilde{G}(\mathfrak{g})$ & $Z(\mathfrak{g})$\\
\hline
$A_r$ & $SU(r+1)$ & $\mathbb{Z}_{r+1}$\\
$B_r$ & $Spin(2r+1)$ & $\mathbb{Z}_2$\\
$C_r$ & $USp(2r)$ & $\mathbb{Z}_2$\\
$D_{2s+1}$ & $Spin(4s+2)$ & $\mathbb{Z}_4$\\
$D_{2s}$ & $Spin(4s)$ & $\mathbb{Z}_2\times\mathbb{Z}_2$\\
$E_6$ & $E_6$ & $\mathbb{Z}_3$\\
$E_7$ & $E_7$ & $\mathbb{Z}_2$\\
$E_8$ & $E_8$ & $1$\\
$F_4$ & $F_4$ & $1$\\
$G_2$ & $G_2$ & $1$\\
\hline
\end{tabular}
\end{center}

% [Tutor Note: 从"$\pi_1$ 是有限生成阿贝尔群"到 (J.2) 之间省掉了两步标准的群论事实：(i) 有限生成阿贝尔群同构于 $\mathbb{Z}^d\oplus$（有限群），而有限群部分对应 $G_1$ 的中心扩张；(ii) 该格的平移必须实现为 $U(1)^d$ 上的作用，故 $\widetilde{G_1}$ 分解为单连通半单部分与 $\mathbb{R}^r$ 的乘积。尤其要注意 $Z$ 必须落在中心里，(J.2) 的商才有群结构——原文只说"$Z$ 是中心的一个子群"，这正是关键。$Z$ 有限这一点决定了 $G_1$ 的整体系数为有限群，而 $\pi_1(G_1)\cong Z$。]

{\small\itshape ［原文标记：还需给出 Dynkin 图以及这些图的自同构群，并解释它与 $G$ 的外自同构之间的关系。］}
% [原文待办标记，原文 p.225；原文以两行星号行围出该段]

{\bfseries 关于可能的扩张，我们能说些什么？}

若我们的群具有正维数，则扩张的存在自动蕴含一个典范同态 $\bar\omega:\pi_0(G)\to\operatorname{Out}(G_1)$。

{\small\itshape ［原文标记：$\ast\ast\ast\ast\ast\ast$］}
% [原文待办标记，原文 p.225；原文此行星号行孤立在段末，含义不明，照录位置]

\subsection{对偶 Coxeter 数} % 原文: J.2 Dual Coxeter Number

单 Lie 代数的对偶 Coxeter 数（dual Coxeter number）由下表给出\footnote{警告：对非单连通（nonsimply laced）的代数还存在一个 Coxeter 数，它与对偶 Coxeter 数不同。}

\begin{center}
\begin{tabular}{|c|c|c|}
\hline
$\mathfrak{g}$ & $\widetilde{G}(\mathfrak{g})$ & $h(\mathfrak{g})$\\
\hline
$A_r$ & $SU(r+1)$ & $r+1$\\
$B_r$ & $Spin(2r+1)$ & $2r-1$\\
$C_r$ & $USp(2r)$ & $r+1$\\
$D_r$ & $Spin(2r)$ & $2r-2$\\
$E_6$ & $E_6$ & $12$\\
$E_7$ & $E_7$ & $18$\\
$E_8$ & $E_8$ & $30$\\
$F_4$ & $F_4$ & $9$\\
$G_2$ & $G_2$ & $4$\\
\hline
\end{tabular}
\end{center}

% [Tutor Note: 这张表有两处极易误用：(i) 表头 $h(\mathfrak{g})$ 是\textbf{对偶} Coxeter 数，而脚注 47 提醒的普通 Coxeter 数在非单连通（nonsimply laced）情形与之不同——$B_r$ 是 $2r$ 对 $2r-1$、$C_r$ 是 $2r$ 对 $r+1$、$F_4$ 是 $12$ 对 $9$、$G_2$ 是 $6$ 对 $4$；而 $A_r$ 与 $D_r$ 两列碰巧相同（$A_r$：$r+1$；$D_r$：$2r-2$），所以拿 $\mathfrak{su}(N)$ 或 $\mathfrak{so}(2r)$ 去试算根本暴露不了混淆，必须用 $B_r$ 或 $C_r$ 才能区分；(ii) $A_r$ 一行即 $h^\vee=r+1=N$，$\mathfrak{su}(N)$ 的对偶 Coxeter 数就是 $N$。]

\section{范畴：群与群胚} % 原文: K. Categories: Groups and Groupoids

\begin{keytakeaways}
  \item 语言：范畴由对象集、hom 集、合成律与单位态射组成；当对象之间存在非平凡自同构时，「相等」不再是可靠的说法，须改用同构与自然变换来表述。
  \item 群的新定位：群恰是只有一个对象、且所有态射都可逆的范畴；群胚（groupoid）则是所有态射都可逆的范畴，其中 $\hom(x,x)$ 是对象 $x$ 的自同构群。
  \item 函子与自然变换：协变与反变只差合成的先后次序（(K.9) 与 (K.10)）；自然变换 $\tau:F_1\to F_2$ 满足 $\tau_Y\circ F_1(f)=F_2(f)\circ\tau_X$，行列式 $\det:GL(n,R)\to U(R)$ 是标准例子。
  \item 等价与物理：全忠实且本质满的函子（(K.36) 与 Freyd 定理）给出范畴等价；按定义，拓扑场论就是配边范畴到 $\mathbf{VECT}$ 的（张量）函子，而规范场论与 orbifold 时空分别对应联络群胚与 $X/\!/G$。
\end{keytakeaways}

范畴（category）这套语言相当抽象，却仍然在弦论与共形场论中找到了晚近的应用。许多物理学家反感范畴语言所牵涉的高度抽象。有些数学家甚至把这一题材称作「抽象废话」（abstract nonsense）。（另一些人则非常认真地对待它。）然而，在弦论与超对称规范场论的进一步形式发展中，它似乎越来越有用。要读懂拓扑场论的任何文献，它也是必需的。我们在此简要说明这套语言的一部分。我们的主要目的是引入关于「群是什么」的一个不同观点，由此导向一个重要的推广：群胚。此外，这一观点也为群上同调（group cohomology）的含义提供了若干非常有趣的洞见。相关的构造在凝聚态物理与拓扑场论中颇为流行。

\textbf{定义} 一个范畴 $\mathcal{C}$ 由下列数据组成：

\begin{enumerate}[label=\alph*.)]
  \item 一个「对象」的集合 $\mathit{Ob}(\mathcal{C})$；
  \item 一族集合 $\hom(X,Y)$，对任意两个对象 $X,Y\in\mathit{Ob}(\mathcal{C})$ 都有定义。$\hom(X,Y)$ 的元素称为「从 $X$ 到 $Y$ 的态射（morphism）」。它们常被写成箭头：
  \begin{equation}
  X\xrightarrow{\ \varphi\ }Y
  \tag{K.1}
  \label{eq:K.1}
  \end{equation}
  \item 一个合成律：
  \begin{equation}
  \hom(X,Y)\times\hom(Y,Z)\to\hom(X,Z)
  \tag{K.2}
  \label{eq:K.2}
  \end{equation}
  \begin{equation}
  (\psi_1,\psi_2)\mapsto\psi_2\circ\psi_1
  \tag{K.3}
  \label{eq:K.3}
  \end{equation}
\end{enumerate}

使得：

\begin{enumerate}
  \item 一个态射 $\varphi$ 唯一决定其源 $X$ 与靶 $Y$。也就是说，对不同的有序对 $(X,Y)$，各 $\hom(X,Y)$ 互不相交。
  \item 对每个 $X\in\mathit{Ob}(\mathcal{C})$ 都存在一个可区分的态射，记作 $1_X\in\hom(X,X)$ 或 $\mathrm{Id}_X\in\hom(X,X)$，它满足
  \begin{equation}
  1_X\circ\varphi=\varphi \qquad \psi\circ 1_X=\psi
  \tag{K.4}
  \label{eq:K.4}
  \end{equation}
  其中 $\varphi\in\hom(Y,X)$、$\psi\in\hom(X,Y)$ 取遍所有 $Y\in\mathit{Ob}(\mathcal{C})$。\footnote{作为练习，请证明这些条件唯一确定了态射 $1_X$。}
  \item 态射的合成是结合的：
  \begin{equation}
  (\psi_1\circ\psi_2)\circ\psi_3=\psi_1\circ(\psi_2\circ\psi_3)
  \tag{K.5}
  \label{eq:K.5}
  \end{equation}
\end{enumerate}

% [Tutor Note: (K.2) 与 (K.3) 的次序约定是读文献时最常见的翻车点：$\hom(X,Y)\times\hom(Y,Z)$ 的第一个变量先作用，故 $(\psi_1,\psi_2)\mapsto\psi_2\circ\psi_1$；这与作者在「注记」第 1 条里刻意区分的 $g\circ f$（视 $f,g$ 为函数）与 $fg$（视其为箭头）两套记法互为对偶。凡在别处见到「$fg$ 表示先 $f$ 后 $g$」的约定，都要先把合成律的次序对齐再代入。另记术语取舍：groupoid 译「群胚」（不取「广群」），functor 译「函子」，monoid 译「幺半群」，而 $\hom$、$\mathit{Ob}$、$\mathit{Mor}$ 等标准符号一律保留、不译。]

另一种有时见到的定义是：范畴由两个集合 $\mathcal{C}_0$（对象）与 $\mathcal{C}_1$（态射）以及两个映射 $p_0:\mathcal{C}_1\to\mathcal{C}_0$、$p_1:\mathcal{C}_1\to\mathcal{C}_0$ 给出。映射 $p_0(f)=x_1\in\mathcal{C}_0$ 是靶映射（range map），$p_1(f)=x_0\in\mathcal{C}_0$ 是源映射（domain map）。在这个替代定义中，范畴由可合成态射集合上的一个合成律给出：
\begin{equation}
\mathcal{C}_2=\{(f,g)\in\mathcal{C}_1\times\mathcal{C}_1\mid p_0(f)=p_1(g)\}
\tag{K.6}
\label{eq:K.6}
\end{equation}
它有时记作 $\mathcal{C}_1{}_{p_1}\!\times_{p_0}\mathcal{C}_1$，称为纤维积（fiber product）。合成律取 $\mathcal{C}_2\to\mathcal{C}_1$，可以画成箭头的合成。若 $f:x_0\to x_1$、$g:x_1\to x_2$，则合成箭头记作 $g\circ f:x_0\to x_2$。该合成律满足下列公理：

\begin{enumerate}
  \item 对每个 $x\in X_0$，$X_1$ 中存在一个单位态射，记作 $1_x$ 或 $\mathrm{Id}_x$，使得对所有可合成的态射 $f,g$ 有 $1_xf=f$ 与 $g1_x=g$。
  \item 合成律是结合的。若 $f,g,h$ 是三个可合成的态射，则 $(hg)f=h(gf)$。
\end{enumerate}

\medskip
\textbf{注记：}

\begin{enumerate}
  \item 定义箭头的合成时，必须做一个重要的记法约定。若 $f:x_0\to x_1$、$g:x_1\to x_2$，则合成箭头是 $x_0\to x_2$ 的箭头。当我们把 $f,g$ 当作函数来想时，我们将写 $g\circ f$；当我们把它们当作箭头来想时，则写 $fg$。
  \item 可以给数据 $X_0,X_1$ 与 $p_0,p_1$ 附加额外的结构，例如拓扑，并要求 $p_0,p_1$ 连续或具有其他性质。
  \item 若存在态射 $\psi\in\hom(Y,X)$ 使得 $\psi\circ\varphi=1_X$ 且 $\varphi\circ\psi=1_Y$，则称态射 $\varphi\in\hom(X,Y)$ 是可逆的（invertible）。若 $X$ 与 $Y$ 之间存在可逆态射，则称 $X$ 与 $Y$ 是同构的对象（isomorphic objects）。使用范畴语言的一个关键理由是：对象可以有非平凡的自同构，即 $\hom(X,X)$ 中可以有除 $1_X$ 之外的可逆元素。当这种情况发生时，谈论对象的「相等」就很棘手，而范畴语言变得非常有用。作为一个具体例子，你不妨想一想下面的问题：「所有 $n$ 维实向量空间都一样吗？」
\end{enumerate}

{\small\itshape ［原文标记：这种对偶记法真是个好主意吗？？］}
% [原文边注/待办标记，原文 p.226；原文以 \ensuremath{\clubsuit}…\ensuremath{\clubsuit} 标在页边，内容是 "Is this dual notation really a good idea??"]

下面是一些简单的范畴例子：

\begin{enumerate}
  \item \textbf{SET}：集合与集合间映射的范畴。\footnote{我们取适当大小的集合族与映射族，以避开集合论中那些恼人的悖论。}
  \item \textbf{TOP}：拓扑空间与连续映射的范畴。
  \item \textbf{TOPH}：拓扑空间与连续映射的同伦类的范畴。
  \item \textbf{MANIFOLD}：流形与适当映射的范畴。我们可以取拓扑流形与流形间的连续映射，也可以取光滑流形与光滑映射作为态射。这两种选择给出两个（非常不同的！）范畴。
  \item \textbf{BORD}(n)：$n$ 维流形的配边范畴。粗略地说，对象是不带边界的 $n$ 维流形，态射是配边。从 $n$ 维流形 $M_1$ 到 $n$ 维流形 $M_2$ 的配边 $Y$ 是一个 $(n+1)$ 维流形，其边界有分解 $\partial Y=(\partial Y)_{in}\amalg(\partial Y)_{out}$，连同微分同胚 $\theta_1:(\partial Y)_{in}\to M_1$ 与 $\theta_2:(\partial Y)_{out}\to M_2$。
  \item \textbf{GROUP}：群与群同态的范畴。注意，若我们把态射取为同构而非同态，就会得到一个非常不同的范畴：其中所有彼此之间态射空间非平凡的对象对，都是同构的群。
  \item \textbf{AB}：阿贝尔群的（子）范畴。
  \item 固定一个群 $G$，令 \textbf{G-SET} 为 $G$-集合的范畴，即带有 $G$-作用的集合 $X$。为简单起见，对 $G$-集合 $X$ 中的点 $x$，我们就把 $G$-作用 $\Phi(g,x)$ 写成 $g\cdot x$。我们要求态射 $f:X_1\to X_2$ 满足 $f(g\cdot x_1)=g\cdot f(x_1)$。
  \item \textbf{VECT}$_\kappa$：域 $\kappa$ 上有限维向量空间的范畴，态射为线性变换。
\end{enumerate}

范畴的一个用处是：它提供了一套语言，能在不同的数学语境中精确地描述「相似的结构」。这样讨论时，为了谈论范畴之间的有趣关系，就必须引入「函子」与「自然变换」的概念。

为了陈述两个范畴之间的关系，我们需要一个「范畴之间的映射」，这就是所谓的函子：

\textbf{定义} 两个范畴 $\mathcal{C}_1$ 与 $\mathcal{C}_2$ 之间的一个函子（functor）由一对映射 $F_{\mathrm{obj}}:\mathit{Obj}(\mathcal{C}_1)\to\mathit{Obj}(\mathcal{C}_2)$ 与 $F_{\mathrm{mor}}:\mathit{Mor}(\mathcal{C}_1)\to\mathit{Mor}(\mathcal{C}_2)$ 组成，使得若
\begin{equation}
x\xrightarrow{\ f\ }y\ \in\hom(x,y)
\tag{K.7}
\label{eq:K.7}
\end{equation}
则
\begin{equation}
F_{\mathrm{obj}}(x)\xrightarrow{\ F_{\mathrm{mor}}(f)\ }F_{\mathrm{obj}}(y)\ \in\hom\big(F_{\mathrm{obj}}(x),F_{\mathrm{obj}}(y)\big)
\tag{K.8}
\label{eq:K.8}
\end{equation}
并且我们还要求 $F_{\mathrm{mor}}$ 与态射的合成相容：这有两种可能的方式。若 $f_1,f_2$ 是可合成的态射，则当
\begin{equation}
F_{\mathrm{mor}}(f_1\circ f_2)=F_{\mathrm{mor}}(f_1)\circ F_{\mathrm{mor}}(f_2)
\tag{K.9}
\label{eq:K.9}
\end{equation}
时，我们说 $F$ 是协变函子（covariant functor）；而当
\begin{equation}
F_{\mathrm{mor}}(f_1\circ f_2)=F_{\mathrm{mor}}(f_2)\circ F_{\mathrm{mor}}(f_1)
\tag{K.10}
\label{eq:K.10}
\end{equation}
时，我们说 $F$ 是反变函子（contravariant functor）。两种情形下我们都还要求\footnote{尽管对所有 $f\in\hom(Y,X)$ 有 $F_{\mathrm{mor}}(\mathrm{Id}_X)\circ F_{\mathrm{mor}}(f)=F_{\mathrm{mor}}(f)$，且对所有 $f\in\hom(X,Y)$ 有 $F_{\mathrm{mor}}(f)\circ F_{\mathrm{mor}}(\mathrm{Id}_X)=F_{\mathrm{mor}}(f)$，但这与「对所有 $\varphi\in\hom(F(Y),F(X))$ 有 $F_{\mathrm{mor}}(\mathrm{Id}_X)\circ\varphi=\varphi$」不是一回事，所以我们需要加上这条额外公理。}
\begin{equation}
F_{\mathrm{mor}}(\mathrm{Id}_X)=\mathrm{Id}_{F(X)}
\tag{K.11}
\label{eq:K.11}
\end{equation}
我们通常省去 $F$ 上的下标，因为由上下文即可看清所指。

\medskip
\textbf{练习}

利用以数据 $p_{0,1}:X_1\to X_0$ 给出的范畴的替代定义，写出函子的概念，并画出相应的交换图。

\medskip
\textbf{练习}（\textit{对偶范畴}）

若 $\mathcal{C}$ 是一个范畴，则对偶范畴 $\mathcal{C}^{opp}$ 就是把所有箭头反向而得到的范畴。更形式地说：对象集合相同，且
\begin{equation}
\hom_{\mathcal{C}^{opp}}(X,Y):=\hom_{\mathcal{C}}(Y,X)
\tag{K.12}
\label{eq:K.12}
\end{equation}
于是对每个态射 $f\in\hom_{\mathcal{C}}(Y,X)$，我们对应一个 $f^{opp}\in\hom_{\mathcal{C}^{opp}}(X,Y)$，使得
\begin{equation}
f_1\circ_{\mathcal{C}^{opp}}f_2=(f_2\circ_{\mathcal{C}}f_1)^{opp}
\tag{K.13}
\label{eq:K.13}
\end{equation}

\begin{enumerate}[label=\alph*.)]
  \item 证明：若 $F:\mathcal{C}\to\mathcal{D}$ 是反变函子，则可以自然地定义一个协变函子 $F:\mathcal{C}^{opp}\to\mathcal{D}$。
  \item 证明：若 $F:\mathcal{C}\to\mathcal{D}$ 是协变函子，则可以自然地定义另一个协变函子 $F^{opp}:\mathcal{C}^{opp}\to\mathcal{D}^{opp}$。
\end{enumerate}

\textbf{例 1：} 每个范畴都有一个到自身的典范函子，称为恒等函子 $\mathrm{Id}_{\mathcal{C}}$。

\textbf{例 2：} 存在一个明显的函子，即从 \textbf{GROUP} 到 \textbf{SET} 的遗忘函子（forgetful functor）。这一想法可推广到许多其他情形：我们「遗忘」掉某些数学结构，而映入一个由更原始的对象构成的范畴。

\textbf{例 3：} 由于 \textbf{AB} 是 \textbf{GROUP} 的子范畴，存在一个明显的函子 $\mathcal{F}:\textbf{AB}\to\textbf{GROUP}$。

\textbf{例 4：} 在下面的一个练习中，你要证明群的阿贝尔化（abelianization）定义了一个函子 $\mathcal{G}:\textbf{GROUP}\to\textbf{AB}$。

\textbf{例 5：} 固定一个群 $G$。那么在上面的讨论中我们已数次用到函子
\begin{equation}
F_G:\textbf{SET}\to\textbf{GROUP}
\tag{K.14}
\label{eq:K.14}
\end{equation}
所用的观察是：若 $X$ 是一个集合，则 $F_G(X)=\mathit{Maps}[X\to G]$ 是一个群。请验证它是一个反变函子：若 $f:X_1\to X_2$ 是集合间的映射，则
\begin{equation}
F_G(X_1)\xleftarrow{\ F_G(f)\ }F_G(X_2)
\tag{K.15}
\label{eq:K.15}
\end{equation}
映射 $F_G(f)$ 通常记作 $f^*$，称为拉回（pull-back）。说得更明确些：若 $\Psi$ 是 $X_2\to G$ 的映射，则 $f^*(\Psi):=\Psi\circ f$ 是 $X_1\to G$ 的映射。

这个函子被用于构造某些非线性 sigma 模型，那是一类靶空间为群 $G$ 的量子场论。把研究对象看作「映入 $G$ 的映射所构成的无穷维群的表示论」这一观点，在 Wess-Zumino-Witten 模型这一特殊情形中极为成功——那是一个具有共形不变性（以及更多性质）的二维量子场论。

\textbf{例 6：} 现在回到范畴 \textbf{G-SET}。固定任意集合 $Y$。那么在上面的讨论中我们已数次用到函子
\begin{equation}
F_{G,Y}:\textbf{G-SET}\to\textbf{G-SET}
\tag{K.16}
\label{eq:K.16}
\end{equation}
所用的观察是：若 $X$ 是 $G$-集合，则 $F_Y(X)=\mathit{Maps}[X\to Y]$ 也是 $G$-集合。为验证它是反变函子，我们写出
\begin{align*}
[g\cdot(f^*\Psi)](x_1)&=(f^*\Psi)(g^{-1}\cdot x_1)\\
&=\Psi\big(f(g^{-1}\cdot x_1)\big)\\
&=\Psi\big(g^{-1}\cdot(f(x_1))\big)\\
&=(g\cdot\Psi)(f(x_1))\tag{K.17}\label{eq:K.17}\\
&=\big(f^*(g\cdot\Psi)\big)(x_1)
\end{align*}
从而 $\Psi\to g\cdot\Psi$ 是 $G$-集合的态射。

这个函子在量子场论中无处不在：若时空具有某种对称性（例如旋转对称或 Poincaré 对称），则同一个群也作用在定义于该时空上的场空间上。

\textbf{例 7：} 固定一个非负整数 $n$ 与一个群 $G$。那么上面讨论过的群上同调（对所有的 $g$ 取平凡扭结 $\omega_g=\mathrm{Id}_A$）定义了一个协变函子
\begin{equation}
H^n(G,\bullet):\textbf{AB}\to\textbf{AB}
\tag{K.18}
\label{eq:K.18}
\end{equation}
为验证它确实是函子，我们需要注意到：若 $\varphi:A_1\to A_2$ 是阿贝尔群的同态，则存在一个诱导同态，通常记作
\begin{equation}
\varphi_*:H^n(G,A_1)\to H^n(G,A_2)
\tag{K.19}
\label{eq:K.19}
\end{equation}
你必须验证 $\mathrm{Id}_*=\mathrm{Id}$ 以及
\begin{equation}
(\varphi_1\circ\varphi_2)_*=(\varphi_1)_*\circ(\varphi_2)_*
\tag{K.20}
\label{eq:K.20}
\end{equation}
严格说来，我们应该把 $\varphi_*$ 记作 $H^n(G,\varphi)$，但对本书作者而言这未免太拘泥了。

\textbf{例 8：} 固定一个非负整数 $n$ 与任意一个群 $A$。那么上面讨论过的群上同调（对所有的 $g$ 取平凡扭结 $\omega_g=\mathrm{Id}_A$）定义了一个反变函子
\begin{equation}
H^n(\bullet,A):\textbf{GROUP}\to\textbf{AB}
\tag{K.21}
\label{eq:K.21}
\end{equation}
为验证它确实是函子，我们需要注意到：若 $\varphi:G_1\to G_2$ 是阿贝尔群的同态，则存在一个诱导同态，通常记作 $\varphi^*$
\begin{equation}
\varphi^*:H^n(G_2,A)\to H^n(G_1,A)
\tag{K.22}
\label{eq:K.22}
\end{equation}
% 原文此处作 "homomorphism of Abelian groups"（用于 $G_1\to G_2$），照录不改。

\textbf{例 9：} \textit{拓扑场论。} 拓扑场论的定义本身就是：它是从流形的配边范畴到向量空间与线性变换范畴的一个函子。关于此点的一些讨论见第 5 章。想了解更多，可参阅若干论文。两个在线资源是：

\begin{quote}
\url{http://www.physics.rutgers.edu/~gmoore/695Fall2015/TopologicalFieldTheory.pdf}（与本讲义非常相似）\\
\url{https://www.ma.utexas.edu/users/dafr/bordism.pdf}
\end{quote}

注意在例 2 中并没有明显的反向函子。当两个范畴之间存在双向的函子时，我们也许会问它们是否在某种意义上是「同一个」。但要说清「同一个」究竟指什么，需要一些小心。

% [Tutor Note: 例 8 处原文写 "If $\varphi:G_1\to G_2$ is a homomorphism of Abelian groups"，但 $G_1,G_2$ 是任意群而非阿贝尔群，这是从例 7 复制句式留下的笔误（应为 group homomorphism）；(K.22) 前的单独一行 $\varphi^*$ 亦系排版残留。凡引用 (K.22) 只取其结论：$H^n(\bullet,A)$ 对群变量反变。]

\textbf{定义} 若 $\mathcal{C}_1$ 与 $\mathcal{C}_2$ 是范畴，而 $F_1:\mathcal{C}_1\to\mathcal{C}_2$ 与 $F_2:\mathcal{C}_1\to\mathcal{C}_2$ 是两个函子，则自然变换（natural transformation）$\tau:F_1\to F_2$ 是这样一条规则：它对每个 $X\in\mathit{Obj}(\mathcal{C}_1)$ 指定一个箭头 $\tau_X:F_1(X)\to F_2(X)$，使得对所有 $X,Y\in\mathit{Obj}(\mathcal{C}_1)$ 与所有 $f\in\hom(X,Y)$ 都有
\begin{equation}
\tau_Y\circ F_1(f)=F_2(f)\circ\tau_X
\tag{K.23}
\label{eq:K.23}
\end{equation}

或者，用图来表示：
\begin{equation}
\begin{array}{ccc}
F_1(X) & \xrightarrow{\ F_1(f)\ } & F_1(Y)\\[2pt]
\downarrow\scriptstyle{\tau_X} & & \downarrow\scriptstyle{\tau_Y}\\[2pt]
F_2(X) & \xrightarrow{\ F_2(f)\ } & F_2(Y)
\end{array}
\tag{K.24}
\label{eq:K.24}
\end{equation}

\textbf{例 1：} \textit{求值映射。} 这是另一个同义反复式的构造，然而仍可能有用。设 $S$ 为任意集合，定义一个函子
\begin{equation}
F_S:\textbf{SET}\to\textbf{SET}
\tag{K.25}
\label{eq:K.25}
\end{equation}
它在对象上的作用为
\begin{equation}
F_S(X):=\mathit{Map}[S\to X]\times S
\tag{K.26}
\label{eq:K.26}
\end{equation}
而若 $\varphi:X_1\to X_2$ 是集合间的映射，则
\begin{equation}
F_S(\varphi):\mathit{Map}[S\to X_1]\times S\to\mathit{Map}[S\to X_2]\times S
\tag{K.27}
\label{eq:K.27}
\end{equation}
由 $F_S(\varphi):(f,s)\mapsto(\varphi\circ f,s)$ 定义。于是我们断言，存在一个到恒等函子的自然变换。对每个集合 $X$ 我们有
\begin{equation}
\tau_X:F_S(X)=\mathit{Map}[S\to X]\times S\to\mathrm{Id}(X)=X
\tag{K.28}
\label{eq:K.28}
\end{equation}
它由 $\tau_X(f,s):=f(s)$ 定义。这被称为「求值映射」（evaluation map）。然后我们需要检验
\begin{equation}
\begin{array}{ccc}
F_S(X) & \xrightarrow{\ \tau_X\ } & X\\[2pt]
\downarrow\scriptstyle{F_S(\varphi)} & & \downarrow\scriptstyle{\varphi}\\[2pt]
F_S(Y) & \xrightarrow{\ \tau_Y\ } & Y
\end{array}
\tag{K.29}
\label{eq:K.29}
\end{equation}
交换。如果把它写开，就会发现这只是一个同义反复。

\textbf{例 2：} \textit{行列式。}\footnote{这个例子用到了一些线性代数的术语，可在下面的「用户手册」第 2 章中找到。} 令 \textbf{COMMRING} 为交换环的范畴，态射为环同态。（于是 $\varphi:R_1\to R_2$ 是阿贝尔群的同态，并且还满足 $\varphi(r\cdot s)=\varphi(r)\cdot\varphi(s)$。）让我们考虑两个函子
\begin{equation}
\textbf{COMMRING}\to\textbf{GROUP}
\tag{K.30}
\label{eq:K.30}
\end{equation}
第一个函子 $F_1$ 把环 $R$ 映到可乘可逆元素构成的乘法群 $U(R)$。这个群常被称为 $R$ 的单位群。若 $\varphi$ 是环的态射且 $r\in U(R_1)$，则 $\varphi(r)\in U(R_2)$，于是由
\begin{equation}
\varphi_*:r\mapsto\varphi(r)
\tag{K.31}
\label{eq:K.31}
\end{equation}
定义的映射 $\varphi_*:U(R_1)\to U(R_2)$ 是群同态。因此 $F_1$ 是函子。第二个函子 $F_2$ 把环 $R$ 映到 $n\times n$ 矩阵的矩阵群 $GL(n,R)$，这些矩阵在 $R$ 中取值且存在逆矩阵。同样，若 $\varphi:R_1\to R_2$ 是态射，则把 $\varphi$ 作用于每个矩阵元素即定义了群同态 $\varphi_*:GL(n,R_1)\to GL(n,R_2)$。现在考虑矩阵 $g\in GL(n,R)$ 的行列式。通常的公式
\begin{equation}
\det(g):=\sum_{\sigma\in S_n}\epsilon(\sigma)\,g_{1,\sigma(1)}\cdot g_{2,\sigma(2)}\cdots g_{n,\sigma(n)}
\tag{K.32}
\label{eq:K.32}
\end{equation}
对 $g\in GL(n,R)$ 完全有意义。此外，
\begin{equation}
\det(g_1g_2)=\det(g_1)\det(g_2)
\tag{K.33}
\label{eq:K.33}
\end{equation}
现在我们断言，行列式定义了自然变换 $\tau:F_1\to F_2$。对每个对象 $R\in\mathit{Ob}(\textbf{COMMRING})$，我们指定态射
\begin{equation}
\tau_R:GL(n,R)\to U(R)
\tag{K.34}
\label{eq:K.34}
\end{equation}
它由 $\tau_R(g):=\det(g)$ 定义。借由 \eqref{eq:K.33}，这确实是范畴 \textbf{GROUP} 中的态射，即它是一个群同态。此外，它满足所需的交换图，因为若 $\varphi:R_1\to R_2$ 是环的态射，则
\begin{equation}
\varphi_*(\det(g))=\det(\varphi_*(g))
\tag{K.35}
\label{eq:K.35}
\end{equation}

% [Tutor Note: 这里方向对不上，是原文的实质笔误：作者声明 $\tau:F_1\to F_2$，但 $F_1(R)=U(R)$、$F_2(R)=GL(n,R)$，而 (K.34) 给出的是 $\tau_R:GL(n,R)\to U(R)$，即 $\tau:F_2\to F_1$。按自然变换的定义 (K.23)，$\tau_X$ 必须从 $F_1(X)$ 出发，所以正确的陈述应是 $\det$ 给出 $F_2\to F_1$ 的自然变换。转写时照原文保留，引用时以 (K.34) 的箭头方向为准，勿照抄 "$\tau:F_1\to F_2$" 一句。]

\textbf{例 3：} \textit{上同调理论中的自然变换。} 如上所强调，上同调群为函子提供了自然的例子。这些不同的上同调群函子之间存在若干有趣的自然变换。

{\small\itshape ［原文标记：我们能不能用它目前发展到的群上同调，讲一个初等的例子呢？？］}
% [原文边注/待办标记，原文 p.232；原文以 \ensuremath{\clubsuit}…\ensuremath{\clubsuit} 标在页边]

\textbf{定义} 若存在函子 $F:\mathcal{C}_1\to\mathcal{C}_2$ 与 $G:\mathcal{C}_2\to\mathcal{C}_1$，连同（通过自然变换给出的）同构 $FG\cong\mathrm{Id}_{\mathcal{C}_2}$ 与 $GF\cong\mathrm{Id}_{\mathcal{C}_1}$，则称两个范畴是等价的（equivalent）。（注意 $FG$ 与 $\mathrm{Id}_{\mathcal{C}_2}$ 都是函子范畴 $\mathrm{FUNCT}(\mathcal{C}_2,\mathcal{C}_2)$ 中的对象，所以说它们同构是有意义的。）

{\small\itshape ［原文标记：应该补一个例子，说明域上有限维向量空间的范畴等价于非负整数的范畴。］}
% [原文边注/待办标记，原文 p.233；原文中 "catetgory" 系拼写错误，照原意译出]

数学中许多重要定理都可以用一句优雅而简洁的话来表述：两个看似不同的范畴其实是等价的。下面是一份（非常有选择性的）清单：\footnote{感谢 G. Segal 的有益讨论，这份清单的准备得益于它。}

\textbf{例 1：} 考虑这样一个范畴：对每个非负整数 $n$ 有一个对象，态射空间为域 $\kappa$ 上可逆 $n\times n$ 矩阵构成的 $GL(n,\kappa)$。这些范畴是等价的。这正是在说：有限维向量空间的唯一不变量是它的维数。

\textbf{例 2：} Lie 群与 Lie 代数之间的基本关系是这样一个陈述：把 Lie 群 $G$ 映到它在单位元处的切空间 $T_1G$ 的函子，是「连通且单连通的 Lie 群范畴」与「有限维 Lie 代数范畴」之间的一个等价。该理论中的一个非平凡定理是：从有限维 Lie 代数范畴到连通且单连通 Lie 群范畴的函子存在。直观上，它由对 Lie 代数元素取指数映射给出。

\textbf{例 3：} 覆盖空间理论讲的是一个范畴等价。一方面我们有带基点空间 $(X,x_0)$ 的覆盖的范畴，另一方面我们有带 $\pi_1(X,x_0)$ 作用的拓扑空间的范畴。与此密切相关的 Galois 理论也可以看作一个范畴等价。

\textbf{例 4：} 含单位的交换 $C^*$-代数范畴等价于紧 Hausdorff 拓扑空间的范畴。这就是 Gelfand 定理。

\textbf{例 5：} 与前一个例子类似，代数几何中的一个要点是：域 $\kappa$ 上的交换代数（不含幂零元）范畴与仿射代数簇的范畴之间存在范畴等价。

\textbf{例 6：} Pontryagin 对偶是局部紧阿贝尔群（及连续同态）范畴与自身之间的一个非平凡自等价。

\textbf{例 7：} Pontryagin 对偶的一个推广是紧群范畴与某类线性张量范畴之间的 Tannaka-Krein 对偶。（其想法是：给定一个满足某些条件的抽象张量范畴，可以构造出一个群；而若该张量范畴是某个紧群的表示范畴，则能还原出那个群。）

\textbf{例 8：} Riemann-Hilbert 对应可以看作平坦联络（又称线性微分方程，又称 D-模）范畴与其单值表示（monodromy representation）范畴之间的等价。

{\small\itshape ［原文标记：这里需要多得多解释。］}
% [原文边注/待办标记，原文 p.234；原文以 \ensuremath{\clubsuit}…\ensuremath{\clubsuit} 标在页边]

在物理中，不同物理理论之间的「对偶」陈述有时可以精确地表述为范畴等价。一个重要的例子是镜像对称：它断言 $X$ 上全纯丛的导出范畴与 $X^\vee$ 上 Lagrangians 的 Fukaya 范畴作为（$A_\infty$-）范畴等价。
% 原文此处作 "($A_\infty$)-)"，多出一个右括号，疑为笔误，照录原样。
但更一般地，弦论与场论中的非平凡对偶对称性都带有强烈的范畴等价的味道。

\medskip
\textbf{练习}（\textit{与自然变换玩耍}）

\begin{enumerate}[label=\alph*.)]
  \item 给定两个范畴 $\mathcal{C}_1,\mathcal{C}_2$，证明自然变换使我们能够定义一个范畴 $\mathrm{FUNCT}(\mathcal{C}_1,\mathcal{C}_2)$，其对象是从 $\mathcal{C}_1$ 到 $\mathcal{C}_2$ 的函子，其态射是自然变换。正因如此，自然变换常被称为「函子的态射」。
  \item 写出恒等函子 $\mathrm{Id}_{\mathcal{C}}$ 到自身的自然变换的含义。证明 $\mathrm{End}(\mathrm{Id}_{\mathcal{C}})$，即恒等函子到自身的所有自然变换之集合，构成一个幺半群（monoid）。
\end{enumerate}

\medskip
\textbf{练习}（\textit{Freyd 定理}）

判断两个范畴是否等价，有一个「实用」的办法，如下所述：

按定义，全忠实函子（fully faithful functor）是这样一个函子 $F:\mathcal{C}_1\to\mathcal{C}_2$：$F_{\mathrm{mor}}$ 在所有 hom 集上都是双射。也就是说，对所有 $X,Y\in\mathit{Obj}(\mathcal{C}_1)$，映射
\begin{equation}
F_{\mathrm{mor}}:\hom(X,Y)\to\hom\big(F_{\mathrm{obj}}(X),F_{\mathrm{obj}}(Y)\big)
\tag{K.36}
\label{eq:K.36}
\end{equation}
是双射。

请证明：$\mathcal{C}_1$ 等价于 $\mathcal{C}_2$ 当且仅当存在全忠实函子 $F:\mathcal{C}_1\to\mathcal{C}_2$，使得 $\mathit{Obj}(\mathcal{C}_2)$ 中任何对象 $\alpha$ 都同构于某个形如 $F(X)$ 的对象（其中 $X\in\mathit{Obj}(\mathcal{C}_1)$）。

\medskip
\textbf{练习}

如上所述，存在一个仅由包含给出的函子 $\textbf{AB}\to\textbf{GROUP}$。

\begin{enumerate}[label=\alph*.)]
  \item 证明阿贝尔化映射 $G\to G/[G,G]$ 定义了函子 $\textbf{GROUP}\to\textbf{AB}$。
  \item 证明非平凡完美群（perfect group）的存在，例如 $A_5$，蕴含这个函子不可能是范畴等价。
\end{enumerate}

除了我们刚才勾勒的这种非常抽象的范畴观之外，非常具体的对象——如群、流形与 orbifold——也可以有益地被看作范畴。

我们总可以把一个范畴画成：对象是点，态射是点之间的有向箭头，如图 39 所示。

作为这一点的极端例子，让我们考虑一个只有一个对象的范畴，但允许有多个态射。对这样的范畴，让我们仔细看看态射集合 $f\in\mathit{Mor}(\mathcal{C})$ 上的结构。我们知道存在一个二元运算，带有一个结合的单位元 $1$。

但这正是幺半群（monoid）的定义！

如果我们还有逆元，就得到一个群。因此：

\textbf{定义} 一个群是只有一个对象、且其所有态射都可逆的范畴。
